% Seccion: Flujo de aire en serpentines
\section{Flujo de aire en serpentines}
\label{sec:cap02_flujo_serpentines}

El flujo de aire a través de serpentines y codos de noventa grados experimenta aceleraciones centrífugas considerables.
Estas fuerzas centrífugas empujan al fluido central de alta velocidad hacia la pared exterior de la curvatura.
Como consecuencia, se generan gradientes de presión transversales que inducen un par de vórtices contrarrotatorios conocidos como vórtices de Dean.
Este movimiento helicoidal secundario actúa como un mezclador térmico natural extremadamente potente.
Los vórtices de Dean aumentan de forma drástica el intercambio de calor convectivo en los laterales del conducto.
Sin embargo, la penalización fluidodinámica se refleja en un incremento del factor de fricción y en pérdidas de carga adicionales.
